\section{Modeling Properties and Containment}
\label{sec:theory}\label{app:containment}
The following statements concern represented feasible sets. They do not imply
that a numerical solver has optimized the complete master or that arbitrary
rounding preserves feasibility.

\subsection{Bounded Products and Epigraphs}
\begin{proposition}\label{prop:mccormick_exact}
For $0\le y,p\le1$, the inequalities $0\le w$, $w\le y$, $w\le p$,
and $w\ge y+p-1$ describe the convex hull of the graph $w=yp$.
If either factor is integral, these inequalities force $w=yp$.
\end{proposition}
\begin{proof}
The lower envelope is $\max(0,y+p-1)$ and the upper envelope is
$\min(y,p)$~\citep{mccormick1976computability}. For $y=0$ they both force
$w=0$; for $y=1$ they force $w=p$. The cases for integral $p$ are symmetric.
\end{proof}
This is an individual-product property, not a convex-hull theorem for the
complete program. For example,
\begin{equation}
y_1=y_2=p=\tfrac12,\quad w_1=w_2=\tfrac12
\label{eq:rlt_counterexample}
\end{equation}
satisfies the separate envelopes and $y_1+y_2=1$, but violates the identity
$w_1+w_2=p$ that holds for exact products. Linking rows
\begin{equation}
\sum_c w_{ick}=p_{\mathrm{assoc}(i,k)}
\label{eq:rlt_linking}
\end{equation}
can strengthen a branch when its intended integral semantics justify them.
They must not be added to an immutable compatibility branch while claiming
containment of every original fractional point.

\begin{corollary}[Disaggregated service epigraphs]\label{cor:v3_hull}
For fixed resource costs $a_c,b_c$, the one-choice epigraph has the perspective
form $\sum_c y_c=1$, $y_c\ge0$, $t_c\ge a_cy_c$, $t_c\ge b_cy_c$,
$T\ge\sum_ct_c$. It gives $T\ge\sum_c y_c\max(a_c,b_c)$ and dominates
$T\ge\max(\sum_ca_cy_c,\sum_cb_cy_c)$.
\end{corollary}
\begin{proof}
For each $c$, minimizing $t_c$ gives $y_c\max(a_c,b_c)$. Summing maxima is
at least either sum separately. With a one-hot choice the two formulations
coincide. For shared variable hardware, every additional coupling must also
be disaggregated consistently; separate product envelopes alone do not prove
the convex hull of the entire HW/SW problem.
\end{proof}

\subsection{An Objective-Preserving Embedding}
\begin{definition}\label{def:seq_scope}
An included schedule branch is a complete, identified program with its own
workload, cost coefficients, hardware and boundary convention. A sequential
control fixes its tile choices and other schedule decisions within that same
contract. A method using different coefficients or unsupported hardware is
not included merely because it has a similar pipeline.
\end{definition}
Write each included branch as
$P_b=\{x:E_bx=e_b,\ G_bx\le g_b,\ \ell_b\le x\le u_b\}$.
The master homogenizes all its rows and finite boxes with $\lambda_b$.
\begin{theorem}[Containment of every included branch]\label{thm:containment}
For every included branch $p$ and every $x\in P_p$, the assignment
$\lambda_p=1$, $x_p=x$, and $\lambda_b=0,x_b=0$ for $b\ne p$ is master-feasible,
with the same objective and hardware readout.
\end{theorem}
\begin{proof}
Every active-branch row and bound becomes its original constraint. Every
inactive-branch row has zero left and right sides; zero satisfies each
homogenized finite box. The branch weights sum to one, active hardware
coordinates are unchanged, and $\sum_b c_b^Tx_b=c_p^Tx$. The argument applies to fractional
points as well as integral schedules.
\end{proof}
\begin{corollary}\label{cor:optimum}
$\inf_{z\in P_M} c_M^Tz\le\min_b\inf_{x\in P_b}c_b^Tx$.
A fixed-tile sequential control represented within an included branch therefore
has a master-feasible embedding at the same modeled cost.
\end{corollary}
\begin{remark}\label{rem:scope}
Containment is not automatic reproduction of every published chip: the actual
branch and its complete contract must be imported. It also does not establish
that the causal block alone contains the old fractional schedule block, or
that a recovery heuristic returns a no-worse integral design. Retaining the
old admitted execution as an eligible incumbent supplies the latter fallback
under matched semantics. Unbounded branches can retain recession directions
at zero weight; we claim the embedding above, not an unconditional
closed-convex-hull identity.
\end{remark}

\subsection{Structure and Safe Pruning}\label{app:invar}
\begin{proposition}\label{prop:invariance}
For fixed graph incidence, variable domains, and candidate library, changes
only to numerical costs and bounds preserve the structural row/column pattern.
\end{proposition}
\begin{proof}
The same indexed row templates address the same coordinates. Hardware-dependent
coefficients can change their values; changing the implementation library or
introducing different lifetimes need not preserve the pattern.
\end{proof}
This permits reuse of indexing and warm starts. A new scaling or norm bound
must still be checked when coefficients change. Pareto pruning similarly needs
componentwise dominance in every resource and interface used by the branch;
a two-axis latency plot alone is insufficient when workspace, traffic tiers,
or numerical formats differ.
